\scanpage{87}

\[ \begin{aligned}
 E&=\text{module d'une ``repr. \unsure{noire|naïve|vraie}'' de }\deleted{\unsure{M}}M,\\
 &\hspace{4cm}\text{de }\unsure{\mathfrak g|\mathcal G}\\&\hspace{4cm}\text{de }\mathfrak g'\\
 &\hspace{4cm}\text{de }G\\&\hspace{4cm}\text{de }G'
 \end{aligned} \]

\[ \begin{aligned}\mathfrak g&=\underline{\mathrm{Lie}}\,G\\ \mathfrak g'&=\underline{\mathrm{Lie}}\,G'\end{aligned} \]

1.2. Données équivalentes

a) fibré \deleted{\unsure{projectif}}\inserted{en droites} proj. $P$ sur $X$, avec section marquée $\sigma$

b) Structure d'extension
\[ 0\longrightarrow\Omega\longrightarrow\mathscr L\longrightarrow\underline O_X\longrightarrow0 \]
$\Omega$ inversible

c) Algèbre d'Azumaya de rg 4 $M$, avec idéal ``spécial'' \deleted{\unsure{$\mathfrak a$}} \unsure{$\mathcal V_\ell$|$\mathcal V_L$|$\mathfrak a$}

d) $\mathfrak g$ avec ss.-Alg. de Borel $\mathscr L$\\
d') $\mathfrak g'$ \rule{4cm}{0.3pt} $\mathscr L'$\\
e) $G$ \rule{2cm}{0.3pt} Gr. \rule{1.5cm}{0.3pt} $B$\\
e') $G'$ \rule{2cm}{0.3pt} gr \rule{1.5cm}{0.3pt} $B'$

% slanted
\margin{\unsure{identifier|vérifier} ! \uline{NB} $\mathfrak g/\Omega\simeq\check{\mathscr L}$ (par Killing) $\simeq\mathscr L\otimes\det\mathscr L^{-1}\simeq\mathscr L\otimes\Omega^{-1}$ avec filtration indiquée}

\deleted{f) Une $\underline O_X$-Algèbre augmentée $B_2$ \deleted{\unsure{trinquée}} \unsure{d'algèbre|d'échelon} \deleted{\unsure{d'idéaux}} \inserted{tronquée à l'}2, de ``rang'' 1}

\inserted{à a) est associé $I=\underline L$} faisceau inv. de degré relatif 1 \uline{rigidifié} par $\sigma^*(\underline L)\simeq\underline O_X$. Alors

\[ \begin{aligned}
 \bullet\quad\mathscr L&=f_*(\underline L),\quad\text{avec structure d'ex. évidente}\\
 &\simeq f_*(\underline L/\mathcal I\underline L)\qquad [\mathcal I=\text{\uline{Idéal}}_{\,\unsure{\sigma}}\text{ de }\sigma(X)\text{ dans }P]\\
 &\simeq\sigma^*(P^1(\underline L))
 \end{aligned} \]

$P=P(\mathscr L)$ avec section provenant du quotient inv. $\underline O_X$ de $\mathscr L$ $={}$ \unsure{cptfion|chpfion} projective du torseur sous $\unsure{\mathbb V|\mathbb W}(\Omega)$ défini par l'extension \inserted{de $\mathscr L$} $={}$ $\unsure{f_{\mathrm{nd}}|\mathrm{End}}\,\mathscr L$, avec idéal augmentation de $\Omega$

% slanted
\margin{(\unsure{adjoint} \deleted{du} prolongé)}

$\mathfrak g=M/\underline O_X\simeq\mathrm{Sym}^2(\mathscr L)\otimes\Omega^{-1}$ ($\underline{\mathrm{NB}}$ $\det\mathscr L\simeq\Omega$) Or $\mathrm{Sym}^2\mathscr L$ est filtré par (décr.) $\underline O_X,\Omega,\Omega^2$, donc $\mathrm{Sym}^2\mathscr L\otimes\Omega^{-1}$ filtré par (croiss.)

\diagram{$\begin{array}{ccc}
 \underbrace{\Omega}&\underline O_X&\Omega^{-1}\\
 \Vert&\Vert&\Vert\\
 \mathrm{Fil}_1&\mathrm{Fil}_2/\mathrm{Fil}_1&\mathrm{Fil}_3/\mathrm{Fil}_2.\\
 \Vert&\Vert&\Vert\\
 \mathrm{gr}_1&\mathrm{gr}_2&\mathrm{gr}_3
 \end{array}$}

% slanted
\margin{expliquer aussi \unsure{comment}, partant de $(\mathfrak g,\mathscr L)$, on trouve $\Omega\hookrightarrow\mathscr L\longrightarrow\underline O_X\longrightarrow0$ $(\Omega=\mathscr L_{\unsure{u|n}},\ldots)$}

\[ \begin{aligned}
 \text{où}\quad \mathrm{Fil}_2&=\Omega^{-1}\otimes\underbrace{(\mathrm{Ker}\ \mathrm{Sym}^2\mathscr L\longrightarrow\mathrm{Sym}^2 O)}_{\simeq\Omega\otimes\mathscr L}\\
 &\simeq\mathscr L
 \end{aligned} \]

donc on a
\[ \underset{\Omega}{\underbrace{\Omega}}\subset\underset{\underline O_X}{\underbrace{\mathscr L}}\subset\underset{\Omega^{-1}}{\underbrace{\mathfrak g}} \]
et $\mathscr L$ est une \uline{ss-Alg. de Borel} ($\Omega$ sa partie unipotente)

\margin{\uline{NB} pour forme de Killing, $\Omega$ et $\mathscr L$ orth. l'un de l'autre, si \unsure{2 inv.}}

On n'explicite pas $(\mathfrak g',\mathscr L')$ \unsure{$\overset{2}{\simeq}$} $(\mathfrak g,\mathscr L)$ si 2 inv), pour $(G,B)/(G',B')$ on a $G'=\underline{\mathrm{Sl}}_{\Omega_S}\mathscr L$, $B'\inserted{\subset G'}={}$ stabilisateur de $\Omega$ $G,B$ déduit par passage au quotient par centre $\mu_2$

[ \uline{\unsure{1.3.|1.4.}} Supposons $X/S$, et considérons des \deleted{connexions sur}\inserted{rel. $S$} les \unsure{Objets|objets} du type envisagé dans 1.1, mais moyennant \deleted{\unsure{l'existen}}\inserted{la donnée de} la structure 1.2. Ainsi, on considère des couples $(P,\sigma,c)$

% slanted
\margin{\uline{plus bas}}

\[ \left\{\begin{aligned} &P\quad\text{fibré en droites projectives sur }X\\ &\sigma\quad\text{section}\\&c\quad\text{connexion de }P\text{ rel. à }S\end{aligned}\right. \]

(mais on ne suppose pas $\sigma$ plate) \unsure{Il nous|Il nous y} sera particulièrement commode d'expliciter en termes de $(\mathscr L,\Omega)$ i.e.
\[ 0\longrightarrow\Omega\longrightarrow\mathscr L\longrightarrow\underline O_X\longrightarrow0 \]
(qui donne $(P,\sigma)$, \inserted{où $(\mathfrak g,\mathscr L)$, comme expliqué} et d'une connexion de $\mathfrak g$. \deleted{On note qu'une} ]

% archive sheet number: 86
